\documentclass[10pt,twocolumn]{article}
\usepackage[a4paper, margin=0.65in]{geometry}
\usepackage{iftex}
\ifPDFTeX
  \usepackage[utf8]{inputenc}
  \usepackage[T1]{fontenc}
\else
  \usepackage{fontspec}
  \usepackage[english, bidi=basic, provide=*]{babel}
  \babelprovide[import, onchar=ids fonts]{english}
  \babelfont{rm}{Noto Sans}
\fi
\usepackage{amsmath,amssymb,amsfonts}
\usepackage{graphicx}
\usepackage{booktabs}
\usepackage{multirow}
\usepackage{caption}
\usepackage{cite}
\usepackage{xcolor}
\usepackage{hyperref}

\definecolor{cGray}{HTML}{7F8C8D}
\definecolor{cBlue}{HTML}{2980B9}
\definecolor{cOrange}{HTML}{E67E22}
\definecolor{cRed}{HTML}{C0392B}
\definecolor{cGreen}{HTML}{27AE60}
\definecolor{cPurple}{HTML}{8E44AD}

\renewcommand{\topfraction}{0.95}
\renewcommand{\bottomfraction}{0.95}
\renewcommand{\textfraction}{0.05}
\renewcommand{\floatpagefraction}{0.85}
\renewcommand{\dbltopfraction}{0.95}
\renewcommand{\dblfloatpagefraction}{0.85}
\setcounter{topnumber}{4}
\setcounter{bottomnumber}{4}
\setcounter{totalnumber}{8}
\setcounter{dbltopnumber}{4}

\hypersetup{
    colorlinks=true,
    linkcolor=blue,
    citecolor=blue,
    urlcolor=blue
}

\title{\textbf{Kinematic-Aware Spatio-Temporal Graph Convolutional Network with Visibility-Guarded $L_2$ Bone Normalization for 3D Hand Skeletal Restoration and Gesture Recognition under Optical Tracking Dropouts}}

\author{
  \textbf{Siriporn Noiumkar}$^{1,*}$ \\
  $^{1}$Department of Creative Media, Faculty of Informatics, Mahasarakham University, \\
  Kantarawichai District, Maha Sarakham 44150, Thailand \\
  \texttt{Email: siriporn.p@msu.ac.th}
}
\date{Submitted to \textit{Multimedia Tools and Applications} (Springer)}

\begin{document}
\maketitle

\begin{abstract}
Markerless infrared stereo hand trackers, such as the Ultraleap Leap Motion Controller (LMC) and its second-generation successor (LMC2), enable high-frame-rate 3D skeletal interaction for multimedia systems, virtual reality, and 3D character animation. However, optical hand tracking across depth and stereo-infrared sensors suffers from two persistent physical limitations: (i) part-level finger self-occlusion and whole-hand temporal burst dropouts caused by line-of-sight obstruction, and (ii) anatomical bone-ratio and camera-distance variations that induce scale-based classification shortcuts and unnatural skeletal distortions during 3D motion retargeting. This paper presents a unified restoration-to-multimedia framework centered on three core contributions: (1) a \textbf{Visibility-Guarded Wrist-Centered $L_2$ Palm-Bone Normalization} formulated in both offline non-causal (\textbf{System A}) and zero-lookahead causal streaming (\textbf{System B}) variants to prevent division-by-zero explosions and scale shortcuts; (2) a \textbf{Kinematic-Aware Spatio-Temporal Graph Convolutional Network (KAR-STGCN)} integrating a two-stage dilated temporal bridge (25-frame receptive field), anatomical graph message passing, bone-length consistency regularization ($\mathcal{L}_{\text{bone}}$), and dual-stream velocity encoding; and (3) \textbf{KAR-Bench Studio}, an open-source diagnostic and Biovision Hierarchy (\texttt{.bvh}) 3D animation retargeting suite. Evaluated under strict subject-independent protocols across SHREC'17 ($N=28$ subjects), FPHA ($N=6$ subjects), real first-generation Leap Motion streams from LMDHG ($N=21$ subjects), and orthogonal Dual-LMC2 hardware recordings ($N=25$ subjects), KAR-STGCN reduces occluded-joint restoration error ($\text{MPJPE}_{\text{miss}}$) relative to CTR-GCN by \textbf{42.1\%} on SHREC'17 Level 3 ($11.42 \pm 0.48\text{ mm}$), \textbf{40.2\%} on SHREC'17 Level 4 ($15.64 \pm 0.58\text{ mm}$), \textbf{41.7\%} on FPHA Level 4 ($19.88 \pm 0.71\text{ mm}$), and \textbf{45.1\%} on natural Dual-LMC2 dropouts ($10.12 \pm 0.45\text{ mm}$ vs. $18.45\text{ mm}$). Furthermore, KAR-STGCN reduces biomechanical joint-angle violation rates in 3D character retargeting from $31.40\%$ ($1\text{€}$ Filter) to \textbf{1.12\%} and achieves $0.88\text{ ms}$ true end-to-end streaming latency ($0.42\text{ ms}$ GPU forward pass with $0\text{ ms}$ lookahead buffer). Subject-level hypothesis tests and linear mixed-effects models confirm statistically robust gains ($p_{\text{Holm}} < 0.002$).

\vspace{0.6em}
\noindent\textbf{Keywords:} 3D Hand Pose Restoration $\cdot$ Optical Motion Capture $\cdot$ Leap Motion Controller 2 $\cdot$ Spatio-Temporal Graph Convolutional Network $\cdot$ $L_2$ Bone Normalization $\cdot$ 3D Character Animation Retargeting
\end{abstract}

% ============================================================================
\section{Introduction}
% ============================================================================
Three-dimensional (3D) skeletal motion capture (MoCap) and hand tracking serve as a foundational interface for interactive multimedia, virtual reality (VR), sign language recognition, and 3D character animation \cite{bregler2007motion, noiumkar2013use, marin2016hand, rastgoo2025deep}. While multi-camera passive optical MoCap studios provide sub-millimeter spatial accuracy, they require reflective markers, controlled studio lighting, and labor-intensive post-processing to repair occluded marker trajectories \cite{tirakoat2011optimized, noiumkar2013use}. Consequently, compact markerless infrared (IR) stereo sensors—such as the first-generation Ultraleap Leap Motion Controller (LMC) \cite{marin2016hand, boulahia2017dynamic} and the second-generation Leap Motion Controller 2 (LMC2) \cite{gilmartin2024duallmc2, matuszczyk2025lmc2}—alongside depth-based skeletal estimators \cite{de2017shrec, kumar2018position}, have become primary sensing modalities for desktop and head-mounted multimedia systems.

Despite advances in sensor resolution, recent metrological evaluations of the LMC2 \cite{matuszczyk2025lmc2}, dual-sensor LMC2 recordings \cite{gilmartin2024duallmc2}, and optical hand tracking studies \cite{boulahia2017dynamic, kumar2018position, liu2024novel} reveal three fundamental technical bottlenecks across optical hand tracking pipelines:
\begin{enumerate}
    \item \textbf{Part-Level Self-Occlusion vs. Whole-Hand Burst Dropouts:} Optical hand trackers exhibit modality-specific failure modes. Depth-based pose estimators \cite{de2017shrec, kumar2018position} frequently suffer from part-level fingertip self-occlusion and localized coordinate drift when fingers cross the palm. Conversely, metrological evaluations of the LMC2 by Matuszczyk et al. \cite{matuszczyk2025lmc2} demonstrate palm positional errors of $5.2\text{--}5.3\text{ mm}$ and static repeatability of $0.4\text{--}0.8\text{ mm}$ within the central interaction zone, but show that the tracking SDK exhibits binary whole-hand dropouts (where all 21 joints drop to zero or freeze simultaneously) when the hand reaches oblique viewing angles or workspace boundaries \cite{boulahia2017dynamic, matuszczyk2025lmc2}.
    \item \textbf{Scale Shortcuts and 3D Character Retargeting Distortions:} Deep classifiers trained on unnormalized 3D coordinates often exploit camera-to-hand distance or subject hand size as a spurious shortcut rather than learning intrinsic finger kinematics. Moreover, mapping uncalibrated or naively imputed trajectories onto 3D character rigs violates rigid bone-length invariants, causing unnatural finger stretching and hyperextended joint angles \cite{noiumkar2013use}.
    \item \textbf{Normalization Singularity under Reference-Joint Loss:} Standard scale normalization divides coordinates by an instantaneous reference bone length (e.g., wrist-to-middle-MCP distance $\|\mathbf{p}_{t,9} - \mathbf{p}_{t,0}\|_2$). When reference joints drop out or collapse while other joints retain residual coordinates, naive division by a small constant $\epsilon$ triggers numerical explosion ($\sim 10^7\text{--}10^8\text{ mm}$), destabilizing graph message passing.
\end{enumerate}

To resolve these challenges without conflating offline post-processing and real-time streaming constraints, we organize our scientific contributions into \textbf{three core pillars}:
\begin{itemize}
    \item \textbf{Core Contribution 1 (Visibility-Guarded Anatomical Normalization):} We unify four heterogeneous 3D hand benchmarks (SHREC'17 \cite{de2017shrec}, FPHA \cite{garcia2018first}, LMDHG \cite{boulahia2017dynamic}, and Dual-LMC2 \cite{gilmartin2024duallmc2}) into a canonical 21-joint topology and formulate a visibility-guarded wrist-centered $L_2$ palm-bone normalization in two distinct operational regimes: \textbf{System A (Offline Sequence-Median)} for post-capture animation cleaning and \textbf{System B (Causal Running-Median)} for zero-lookahead streaming. We prove its robustness against reference-joint dropouts, $5\%\text{--}20\%$ visibility mask corruption, and $\times 0.70\text{--}\times 1.40$ scale perturbations.
    \item \textbf{Core Contribution 2 (Kinematic-Aware Spatio-Temporal Restoration — KAR-STGCN):} We design a lightweight network ($0.42\text{M}$ parameters, $0.18\text{ GFLOPs}$) that couples a two-stage dilated temporal bridge (25-frame receptive field) with anatomical graph convolutions, explicit bone-length consistency regularization ($\mathcal{L}_{\text{bone}}$), and dual-stream position-velocity classification.
    \item \textbf{Core Contribution 3 (Unified Restoration-to-Multimedia Application Framework):} We develop \textbf{KAR-Bench Studio}, an open-source multimedia software suite providing real-time bone-length oscilloscopy, end-to-end latency profiling ($0.88\text{ ms}$ total streaming latency), TorchScript export, and automated Biovision Hierarchy (\texttt{.bvh}) 3D character retargeting that reduces biomechanical joint-angle violations to $1.12\%$.
\end{itemize}

% ============================================================================
\section{Related Work}
% ============================================================================
\subsection{3D Hand Pose Estimation and Optical Tracking Metrology}
Early optical MoCap research established that marker placement optimization and post-capture trajectory cleaning are essential for biomechanical fidelity in sports science and character animation \cite{bregler2007motion, tirakoat2011optimized, noiumkar2013use}. For markerless hand tracking, the first-generation LMC \cite{marin2016hand, boulahia2017dynamic} and second-generation LMC2 \cite{gilmartin2024duallmc2, matuszczyk2025lmc2} estimate 3D skeletal joints from stereo IR cameras. Matuszczyk et al. \cite{matuszczyk2025lmc2} recently quantified LMC2 accuracy ($5.2\text{--}5.3\text{ mm}$ palm error, $0.4\text{--}0.8\text{ mm}$ repeatability) and documented that tracking loss occurs as whole-hand dropouts at extreme angles. While Gil-Mart{\'\i}n et al. \cite{gilmartin2024duallmc2} mitigated single-view occlusion hardware-wise using dual orthogonal LMC2 sensors, single-sensor multimedia deployments require algorithmic restoration, and dual-sensor setups provide an ideal cross-view reference for benchmarking single-sensor dropout recovery.

\subsection{Occlusion-Robust Hand Tracking and Time-Series Imputation}
In commercial VR/AR systems, speed-adaptive low-pass filters such as the One-Euro ($1\text{€}$) Filter \cite{casiez2012oneeuro} and Rauch-Tung-Striebel (RTS) Kalman smoothers are widely used to suppress high-frequency jitter, but they rely on zero-order hold or linear extrapolation during multi-frame tracking loss. In deep multivariate time-series imputation, bidirectional recurrent models (BRITS \cite{cao2018brits}) and diagonally masked self-attention networks (SAITS \cite{du2023saits}) achieve strong point-wise imputation. However, generic imputers ignore 3D skeletal kinematic chains, causing severe bone-length shrinkage or elongation during rapid finger articulation.

\subsection{Skeleton-Based Action and Gesture Recognition}
In \textit{Multimedia Tools and Applications} and computer vision venues, 3D hand gesture recognition has evolved from geometric features \cite{marin2016hand, boulahia2017dynamic, kumar2018position} and fast kinematic 1D-CNNs such as DD-Net \cite{yang2019ddnet} to Spatio-Temporal Graph Convolutional Networks (ST-GCN \cite{yan2018spatial}, 2s-AGCN \cite{shi2019twostream}, MS-G3D \cite{liu2020msg3d}, CTR-GCN \cite{chen2021ctrgcn}) and recent topology-aware architectures such as BlockGCN \cite{zhou2024blockgcn} and residual attention networks \cite{asteasuain2022applying, liu2024novel, rastgoo2025deep, dhanjal2025multilingual}. Because standard GCNs assume fully observed node coordinates, feeding zero-dropped frames directly into graph adjacency layers propagates corrupted features to neighboring joints.

\subsection{Motion Capture Reconstruction and Anatomical Constraints}
Mapping raw or naively smoothed skeletal streams onto 3D character rigs frequently produces visually jarring limb stretching and invalid joint rotations \cite{tirakoat2011optimized, noiumkar2013use}. Unlike disjoint two-stage pipelines that chain a generic time-series imputer with a separate classifier, KAR-STGCN jointly optimizes coordinate imputation, bone-length preservation, and velocity-informed gesture recognition within a single end-to-end architecture.

% ============================================================================
\section{Proposed Methodology}
% ============================================================================
\subsection{Canonical 21-Joint Topology Alignment}
We unify four heterogeneous 3D hand datasets into a canonical 21-joint hand graph $\mathcal{G} = (\mathcal{V}, \mathcal{E})$ with $|\mathcal{V}|=21$ joints and $|\mathcal{E}|=20$ anatomical bone segments: wrist ($0$), thumb ($1\text{--}4$), index ($5\text{--}8$), middle ($9\text{--}12$), ring ($13\text{--}16$), and pinky ($17\text{--}20$). For \textbf{SHREC'17} \cite{de2017shrec} ($J=22$), the redundant palm center (index 1) is removed. For \textbf{FPHA} \cite{garcia2018first} ($J=21$ magnetic sensors), the finger-root-first indexing is bijectively permuted to $\mathcal{V}$. For \textbf{LMDHG} \cite{boulahia2017dynamic} (first-generation Leap Motion streams) and \textbf{Dual-LMC2} \cite{gilmartin2024duallmc2} (second-generation Ultraleap LMC2 streams), we extract the active dominant hand and prune auxiliary forearm/palm-center nodes to retain the exact 21 hand joints matching $\mathcal{V}$. All sequences are resampled to $T=60$ frames ($\mathbf{X}_{\text{mm}} \in \mathbb{R}^{T \times 21 \times 3}$).

\subsection{Visibility-Guarded Wrist-Centered $L_2$ Normalization: Offline vs. Causal Regimes}
Let $\mathbf{p}_{t,j} \in \mathbb{R}^3$ denote the 3D position (in mm) of joint $j$ at frame $t$, and let $v_{t,j} \in \{0,1\}$ denote its binary visibility indicator. When explicit per-joint confidence is not provided by the sensor SDK, we infer $\hat{v}_{t,j}$ causally from frame-to-frame kinematics: $\hat{v}_{t,j}=0$ if $\|\mathbf{p}_{t,j}\|_2 < 10^{-3}\text{ mm}$ (zero-collapse), $\|\mathbf{p}_{t,j} - \mathbf{p}_{t-1,j}\|_2 < 10^{-5}\text{ mm}$ for $\ge 3$ consecutive frames (tracking freeze), or $\|\mathbf{p}_{t,j} - \mathbf{p}_{t-1,j}\|_2 > 60\text{ mm/frame}$ (unphysical jump).

To prevent division-by-zero explosions when the wrist ($0$) or middle MCP joint ($9$) drops out, we define the set of valid reference frames as $\mathcal{T}_{\text{valid}} = \{t \in [1, T] \mid v_{t,0}=1 \land v_{t,9}=1 \land \|\mathbf{p}_{t,9}-\mathbf{p}_{t,0}\|_2 > \tau_{\min}\}$ with $\tau_{\min} = 15.0\text{ mm}$ (since adult human palm length is $70\text{--}100\text{ mm}$). Let $t^* = \min\{k \mid v_{k,0}=1\}$ denote the first valid wrist frame; the causal last-observed valid wrist anchor $\mathbf{p}_{t,0}^{\text{valid}}$ is defined recursively as:
\begin{equation}
    \mathbf{p}_{t,0}^{\text{valid}} = \begin{cases} 
    \mathbf{p}_{t,0}, & \text{if } v_{t,0} = 1, \\
    \mathbf{p}_{t-1,0}^{\text{valid}}, & \text{if } v_{t,0} = 0 \text{ and } t > 1, \\
    \mathbf{p}_{t^*,0}, & \text{otherwise.}
    \end{cases}
    \label{eq:valid_wrist}
\end{equation}
To strictly separate offline post-processing from real-time streaming (addressing causal dependency requirements), we define two distinct normalization operators:
\begin{itemize}
    \item \textbf{System A (Offline Non-Causal Normalization):} Uses the full sequence median $d_{\text{ref}}^{\text{off}} = \operatorname{median}_{k \in \mathcal{T}_{\text{valid}}} \|\mathbf{p}_{k,9} - \mathbf{p}_{k,0}\|_2$.
    \item \textbf{System B (Causal Streaming Normalization):} Uses only past and current frames $\mathcal{T}_{\text{valid}}^{\le t} = \{k \le t \mid k \in \mathcal{T}_{\text{valid}}\}$ to compute the running median $d_{\text{ref}}^{\text{causal}}(t) = \operatorname{median}_{k \in \mathcal{T}_{\text{valid}}^{\le t}} \|\mathbf{p}_{k,9} - \mathbf{p}_{k,0}\|_2$ with zero future-frame leakage.
\end{itemize}
In both systems, the normalized coordinate $\tilde{\mathbf{p}}_{t,j}$ is computed as:
\begin{equation}
    \tilde{\mathbf{p}}_{t,j} = v_{t,j} \cdot \frac{\mathbf{p}_{t,j} - \mathbf{p}_{t,0}^{\text{valid}}}{d_{t,\text{palm}}}, \quad d_{t,\text{palm}} = \begin{cases} \|\mathbf{p}_{t,9}-\mathbf{p}_{t,0}\|_2, & t \in \mathcal{T}_{\text{valid}}, \\ d_{\text{ref}}, & t \notin \mathcal{T}_{\text{valid}}, \end{cases}
    \label{eq:l2norm}
\end{equation}
defaulting to $80.0\text{ mm}$ prior to the first valid frame.

\begin{figure*}[!t]
\centering
\setlength{\unitlength}{1mm}
\begin{picture}(172, 62)
  % Legend at top
  \put(8, 57){\textcolor{cGray}{\rule{3.5mm}{2.2mm}} \scriptsize Kinematic-Filter}
  \put(40, 57){\textcolor{cPurple}{\rule{3.5mm}{2.2mm}} \scriptsize SAITS (Imputation)}
  \put(76, 57){\textcolor{cBlue}{\rule{3.5mm}{2.2mm}} \scriptsize ST-GCN}
  \put(102, 57){\textcolor{cOrange}{\rule{3.5mm}{2.2mm}} \scriptsize CTR-GCN}
  \put(132, 57){\textcolor{cRed}{\rule{3.5mm}{2.2mm}} \scriptsize \textbf{Proposed KAR-STGCN}}

  % ---------------- Panel (a): SHREC'17 ----------------
  \put(12, 51){\small\textbf{(a) SHREC'17: Overall MPJPE (mm) $\downarrow$}}
  \put(10, 10){\line(1,0){72}}
  \put(10, 10){\line(0,1){38}}
  \put(5, 10){\scriptsize 0}
  \put(4, 19.5){\scriptsize 10}
  \put(4, 29.0){\scriptsize 20}
  \put(4, 38.5){\scriptsize 30}

  % Level 1 bars
  \put(13, 10){\textcolor{cGray}{\rule{2.3mm}{3.9mm}}}
  \put(15.6, 10){\textcolor{cPurple}{\rule{2.3mm}{2.5mm}}}
  \put(18.2, 10){\textcolor{cBlue}{\rule{2.3mm}{2.3mm}}}
  \put(20.8, 10){\textcolor{cOrange}{\rule{2.3mm}{2.0mm}}}
  \put(23.4, 10){\textcolor{cRed}{\rule{2.3mm}{1.7mm}}}
  \put(14, 5){\scriptsize Level 1}

  % Level 2 bars
  \put(30, 10){\textcolor{cGray}{\rule{2.3mm}{11.2mm}}}
  \put(32.6, 10){\textcolor{cPurple}{\rule{2.3mm}{8.0mm}}}
  \put(35.2, 10){\textcolor{cBlue}{\rule{2.3mm}{7.5mm}}}
  \put(37.8, 10){\textcolor{cOrange}{\rule{2.3mm}{6.6mm}}}
  \put(40.4, 10){\textcolor{cRed}{\rule{2.3mm}{4.6mm}}}
  \put(31, 5){\scriptsize Level 2}

  % Level 3 bars
  \put(47, 10){\textcolor{cGray}{\rule{2.3mm}{23.2mm}}}
  \put(49.6, 10){\textcolor{cPurple}{\rule{2.3mm}{13.5mm}}}
  \put(52.2, 10){\textcolor{cBlue}{\rule{2.3mm}{14.2mm}}}
  \put(54.8, 10){\textcolor{cOrange}{\rule{2.3mm}{12.3mm}}}
  \put(57.4, 10){\textcolor{cRed}{\rule{2.3mm}{8.2mm}}}
  \put(48, 5){\scriptsize Level 3}

  % Level 4 bars
  \put(64, 10){\textcolor{cGray}{\rule{2.3mm}{31.2mm}}}
  \put(66.6, 10){\textcolor{cPurple}{\rule{2.3mm}{18.4mm}}}
  \put(69.2, 10){\textcolor{cBlue}{\rule{2.3mm}{19.6mm}}}
  \put(71.8, 10){\textcolor{cOrange}{\rule{2.3mm}{16.9mm}}}
  \put(74.4, 10){\textcolor{cRed}{\rule{2.3mm}{10.8mm}}}
  \put(65, 5){\scriptsize Level 4}

  % ---------------- Panel (b): FPHA ----------------
  \put(98, 51){\small\textbf{(b) FPHA: Overall MPJPE (mm) $\downarrow$}}
  \put(96, 10){\line(1,0){72}}
  \put(96, 10){\line(0,1){38}}
  \put(91, 10){\scriptsize 0}
  \put(89, 17.6){\scriptsize 10}
  \put(89, 25.2){\scriptsize 20}
  \put(89, 32.8){\scriptsize 30}
  \put(89, 40.4){\scriptsize 40}

  % Level 1 bars
  \put(99, 10){\textcolor{cGray}{\rule{2.3mm}{3.4mm}}}
  \put(101.6, 10){\textcolor{cPurple}{\rule{2.3mm}{2.7mm}}}
  \put(104.2, 10){\textcolor{cBlue}{\rule{2.3mm}{2.5mm}}}
  \put(106.8, 10){\textcolor{cOrange}{\rule{2.3mm}{2.3mm}}}
  \put(109.4, 10){\textcolor{cRed}{\rule{2.3mm}{1.9mm}}}
  \put(100, 5){\scriptsize Level 1}

  % Level 2 bars
  \put(116, 10){\textcolor{cGray}{\rule{2.3mm}{11.1mm}}}
  \put(118.6, 10){\textcolor{cPurple}{\rule{2.3mm}{8.7mm}}}
  \put(121.2, 10){\textcolor{cBlue}{\rule{2.3mm}{8.3mm}}}
  \put(123.8, 10){\textcolor{cOrange}{\rule{2.3mm}{7.4mm}}}
  \put(126.4, 10){\textcolor{cRed}{\rule{2.3mm}{5.2mm}}}
  \put(117, 5){\scriptsize Level 2}

  % Level 3 bars
  \put(133, 10){\textcolor{cGray}{\rule{2.3mm}{24.2mm}}}
  \put(135.6, 10){\textcolor{cPurple}{\rule{2.3mm}{14.9mm}}}
  \put(138.2, 10){\textcolor{cBlue}{\rule{2.3mm}{15.8mm}}}
  \put(140.8, 10){\textcolor{cOrange}{\rule{2.3mm}{13.8mm}}}
  \put(143.4, 10){\textcolor{cRed}{\rule{2.3mm}{9.4mm}}}
  \put(134, 5){\scriptsize Level 3}

  % Level 4 bars
  \put(150, 10){\textcolor{cGray}{\rule{2.3mm}{34.5mm}}}
  \put(152.6, 10){\textcolor{cPurple}{\rule{2.3mm}{21.2mm}}}
  \put(155.2, 10){\textcolor{cBlue}{\rule{2.3mm}{22.4mm}}}
  \put(157.8, 10){\textcolor{cOrange}{\rule{2.3mm}{19.5mm}}}
  \put(160.4, 10){\textcolor{cRed}{\rule{2.3mm}{12.5mm}}}
  \put(151, 5){\scriptsize Level 4}
\end{picture}
\caption{Subject-independent comparison of overall 3D skeletal restoration error ($\text{MPJPE}_{\text{all}}$ in mm $\downarrow$) across all four optical degradation levels on (a) SHREC'17 ($N=28$) and (b) FPHA ($N=6$) benchmarks.}
\label{fig:mpjpe_bars}
\end{figure*}

\subsection{KAR-STGCN Architecture: System A (Offline) vs. System B (Causal)}
Given normalized coordinates $\tilde{\mathbf{X}} \in \mathbb{R}^{T \times 21 \times 3}$ and visibility mask $\mathbf{V} \in \{0,1\}^{T \times 21}$, we apply a two-stage dilated depthwise temporal bridge $\mathcal{B}_{\text{temp}}$ consisting of two cascaded 1D convolutions of kernel size $K_b=9$ with dilation factors $d_1=1$ and $d_2=2$:
\begin{itemize}
    \item In \textbf{System A (Offline Restoration)}, symmetric padding yields a bidirectional receptive field of $25$ frames ($\pm 12$ frames, requiring a $12$-frame $/ 100.0\text{ ms}$ lookahead buffer at $120\text{ Hz}$).
    \item In \textbf{System B (Causal Streaming)}, left-only causal padding restricts convolutions strictly to $\{t-24, \dots, t\}$, achieving a $25$-frame causal history with \textbf{zero lookahead delay ($0.0\text{ ms}$ buffer)}.
\end{itemize}
The bridge imputes occluded coordinates via visibility gating:
\begin{equation}
    \mathbf{X}_{\text{imp}} = \mathbf{V} \odot \tilde{\mathbf{X}} + (1 - \mathbf{V}) \odot \mathcal{B}_{\text{temp}}(\tilde{\mathbf{X}}).
    \label{eq:bridge}
\end{equation}
The concatenated tensor $\mathbf{F}_{\text{in}} = [\mathbf{X}_{\text{imp}} \,\|\, \mathbf{V}] \in \mathbb{R}^{T \times 21 \times 4}$ is projected to $C=64$ channels and processed through normalized skeletal graph convolutions $\hat{\mathbf{A}} = \mathbf{D}^{-1/2}(\mathbf{A}+\mathbf{I})\mathbf{D}^{-1/2}$ \cite{yan2018spatial} and multi-scale temporal convolutions ($3\times 1$ and $5\times 1$). The restored pose $\hat{\mathbf{X}}$ is obtained via a visibility-weighted residual gate ($\alpha_{\text{vis}}=0.15$):
\begin{equation}
    \hat{\mathbf{X}} = \mathbf{X}_{\text{imp}} + \left(\alpha_{\text{vis}}\mathbf{V} + (1-\alpha_{\text{vis}})(1-\mathbf{V})\right) \odot \phi_{\text{dec}}(\mathbf{H}^{(L)}).
    \label{eq:gate}
\end{equation}
For gesture classification, restored positions $\hat{\mathbf{X}}$ and first-order temporal velocities $\Delta\hat{\mathbf{X}}_t = \hat{\mathbf{X}}_t - \hat{\mathbf{X}}_{t-1}$ feed a batch-normalized dilated 1D convolutional head. The joint training objective combines Smooth-$L_1$ coordinate loss, Bone-Length Consistency Error ($\mathcal{L}_{\text{bone}}$), and class-balanced cross-entropy ($\mathcal{L}_{\text{cls}}$):
\begin{equation}
\begin{split}
    \mathcal{L}_{\text{bone}} &= \frac{1}{T|\mathcal{E}|}\sum_{t=1}^T \sum_{(u,v)\in\mathcal{E}} \Big| \|\hat{\mathbf{p}}_{t,u}-\hat{\mathbf{p}}_{t,v}\|_2 \\
    &\quad - \|\mathbf{p}^*_{t,u}-\mathbf{p}^*_{t,v}\|_2 \Big|,
\end{split}
\label{eq:lbone}
\end{equation}
\begin{equation}
\begin{split}
    \mathcal{L}_{\text{total}} &= \mathcal{L}_{\text{SmoothL1}}(\hat{\mathbf{X}}, \mathbf{X}^*) + \lambda_{\text{bone}} \mathcal{L}_{\text{bone}} \\
    &\quad + \lambda_{\text{cls}} \mathcal{L}_{\text{cls}}(\hat{\mathbf{y}}, y),
\end{split}
\label{eq:ltotal}
\end{equation}
where $\lambda_{\text{bone}}=0.60$ and $\lambda_{\text{cls}}=0.50$.

% ============================================================================
\section{Experimental Protocol}
% ============================================================================
\subsection{Datasets and True Subject-Level Splits}
All evaluations strictly enforce \textbf{Subject-Independent Cross-Validation} where no subject appears in both training and test splits:
\begin{enumerate}
    \item \textbf{SHREC'17} \cite{de2017shrec}: $2,800$ sequences across $N=28$ subjects (14-class and 28-class protocols). Unit of statistical analysis: $N=28$ independent subjects ($df=27$).
    \item \textbf{FPHA} \cite{garcia2018first}: $1,175$ sequences across $N=6$ subjects and $45$ object-manipulation action classes with true 6-DoF magnetic sensor ground truth. To prevent pseudoreplication across random seeds, subject-level paired tests strictly use $N=6$ independent subjects ($df=5$, averaging across 5 seeds per subject), supplemented by a hierarchical Linear Mixed-Effects Model (LMM) with random intercepts for subject and seed ($6 \times 5 = 30$ runs).
    \item \textbf{LMDHG (Real First-Generation Leap Motion)} \cite{boulahia2017dynamic}: $608$ gesture instances across $N=21$ subjects ($df=20$) and $13$ dynamic gesture classes. We separate LMDHG evaluation into: (i) \textbf{Natural Unmasked Recognition} (Macro-F1 on raw recorded Leap Motion streams), and (ii) \textbf{Held-Out Contiguous Masking Restoration} ($\text{MPJPE}_{\text{miss}}$ and BLCE on verified segments masked with $5\text{--}13$ frame whole-hand bursts).
    \item \textbf{Orthogonal Dual-LMC2 (Real Second-Generation LMC2 with Cross-View Ground Truth)} \cite{gilmartin2024duallmc2}: $375$ sequences across $N=25$ subjects ($df=24$) and $15$ dynamic hand gestures recorded simultaneously by two orthogonal Ultraleap LMC2 sensors. Primary Sensor 1 (horizontal desktop placement) undergoes natural self-occlusion and oblique whole-hand tracking dropouts, while Orthogonal Sensor 2 (registered into Sensor 1's coordinate system via Kabsch $\text{SE}(3)$ rigid alignment over high-confidence calibration frames) provides an independent physical ground-truth reference during Sensor 1 tracking loss.
\end{enumerate}

\subsection{Six-Type Optical Dropout Simulator and Progressive Levels}
To systematically model both depth-camera part occlusions \cite{de2017shrec} and stereo-IR whole-hand tracking failures \cite{boulahia2017dynamic, matuszczyk2025lmc2}, our simulator incorporates six elementary failure modes: \textbf{Type A} (distal fingertip occlusion, $j \in \{4,8,12,16,20\}$), \textbf{Type B} (partial finger phalanx loss), \textbf{Type C} (instantaneous whole-hand loss), \textbf{Type D} (contiguous whole-hand temporal burst dropout of $5\text{--}40$ frames), \textbf{Type E} (frozen-frame coordinate lock), and \textbf{Type F} (compound failure). These are grouped into four benchmark levels: \textbf{Level 1} (Baseline Jitter, $\sigma=0.8\text{ mm}$), \textbf{Level 2} (Type A+B Part Occlusion, $\sigma=2.2\text{ mm} + 20\%$ distal drop), \textbf{Level 3} (Type C+D+E Whole-Hand Burst Dropout, $\sigma=3.2\text{ mm} + 5\text{--}13$ contiguous frames of complete 21-joint loss), and \textbf{Level 4} (Type F Compound Whole-Hand Burst + $15\%$ Part-Level Drop + $10:1$ Long-Tail Imbalance, $\approx 26.5\%$ overall missing data).

All baselines receive identical corrupted inputs, visibility masks $\mathbf{V}$, train/val/test splits, and individually tuned learning rates ($\text{lr}=1\times 10^{-3}$ with warmup for SAITS \cite{du2023saits}; $2.5\times 10^{-3}$ for GCNs over 40 epochs, batch size 16, across 5 seeds: \texttt{42, 100, 2026, 777, 999}).

% ============================================================================
\section{Experimental Results}
% ============================================================================
\subsection{3D Hand Skeletal Restoration Performance}
Figure \ref{fig:mpjpe_bars} and Table \ref{tab:main_restoration} present the restoration results across SHREC'17 and FPHA. While VR filters ($1\text{€}$ Filter \cite{casiez2012oneeuro}) and Savitzky-Golay smoothing perform adequately under Level 1 jitter ($4.12\text{--}4.48\text{ mm}$), they fail under Level 3 whole-hand bursts ($\text{MPJPE}_{\text{miss}} > 41.8\text{ mm}$). Standard GCNs (ST-GCN, MS-G3D, CTR-GCN, BlockGCN \cite{zhou2024blockgcn}) incur $\text{MPJPE}_{\text{miss}} > 18.6\text{ mm}$ on SHREC'17 Level 3 due to zero-filled feature propagation. Equipping CTR-GCN with SAITS imputation (\textbf{SAITS + CTR-GCN}) or our dilated bridge (\textbf{CTR-GCN + $\mathcal{B}_{\text{temp}}$}) reduces $\text{MPJPE}_{\text{miss}}$ from $19.74\text{ mm}$ to $14.85\text{ mm}$ and $13.60\text{ mm}$, proving that explicit temporal bridging benefits standard GCNs. Nevertheless, \textbf{System A (Offline KAR-STGCN)} achieves the lowest $\text{MPJPE}_{\text{miss}}$ across all benchmarks: \textbf{11.42 $\pm$ 0.48 mm} on SHREC'17 Level 3 (\textbf{42.1\%} reduction vs. CTR-GCN), \textbf{15.64 $\pm$ 0.58 mm} on SHREC'17 Level 4 (\textbf{40.2\%} reduction), \textbf{15.12 $\pm$ 0.55 mm} on FPHA Level 3 (\textbf{42.7\%} reduction), and \textbf{19.88 $\pm$ 0.71 mm} on FPHA Level 4 (\textbf{41.7\%} reduction). Crucially, \textbf{System B (Causal Streaming KAR-STGCN)} achieves $12.35 \pm 0.53\text{ mm}$ on SHREC'17 Level 3 and $16.45 \pm 0.62\text{ mm}$ on Level 4 with \textbf{zero lookahead delay ($0.88\text{ ms}$ total end-to-end latency)}, outperforming offline SAITS+CTR-GCN ($14.85\text{ mm}$, $251.32\text{ ms}$ total latency).

\begin{table*}[!t]
\centering
\caption{Subject-independent 3D skeletal restoration performance ($\text{Mean} \pm \text{SD}$) on SHREC'17 ($N=28$ subjects) and FPHA ($N=6$ subjects). We explicitly report both GPU Forward Latency and Total End-to-End Latency (Lookahead Buffer + CPU Norm/Denorm + GPU Forward Pass at $120\text{ Hz}$).}
\label{tab:main_restoration}
\resizebox{\textwidth}{!}{%
\begin{tabular}{llcccccccc}
\toprule
\textbf{Dataset \& Level} & \textbf{Method / System} & \textbf{Params / FLOPs} & \textbf{$\text{MPJPE}_{\text{all}}$ (mm)} $\downarrow$ & \textbf{$\text{MPJPE}_{\text{miss}}$ (mm)} $\downarrow$ & \textbf{PA-MPJPE (mm)} $\downarrow$ & \textbf{BLCE (mm)} $\downarrow$ & \textbf{Jitter ($\text{mm/f}^2$)} $\downarrow$ & \textbf{GPU Fwd (ms)} & \textbf{End-to-End Lat. (ms)} $\downarrow$ \\
\midrule
\multirow{3}{*}{\shortstack[l]{SHREC'17\\(Level 1 \& 2)}} 
 & Kinematic-Filter (L1 / L2) & -- & $4.12$ / $11.85$ & -- / $24.60$ & $2.85$ / $7.42$ & $2.90$ / $8.12$ & $0.95$ / $2.84$ & -- & $25.08$ (3f buf) \\
 & CTR-GCN \cite{chen2021ctrgcn} (L1 / L2) & 0.54M / 0.23G & $1.94$ / $6.75$ & -- / $13.80$ & $1.32$ / $4.45$ & $1.18$ / $4.10$ & $0.52$ / $1.65$ & $0.51$ & $100.97$ (Offline) \\
 & \textbf{KAR-STGCN (Sys. A, L1 / L2)} & \textbf{0.42M / 0.18G} & $\mathbf{1.78}$ / $\mathbf{4.62}$ & -- / $\mathbf{8.45}$ & $\mathbf{1.15}$ / $\mathbf{3.10}$ & $\mathbf{0.92}$ / $\mathbf{2.35}$ & $\mathbf{0.38}$ / $\mathbf{1.04}$ & $0.42$ & $100.88$ (Offline) \\
\midrule
\multirow{11}{*}{\shortstack[l]{SHREC'17\\(Level 3:\\Whole-Hand\\Burst Drop)}} 
 & $1\text{€}$ Filter \cite{casiez2012oneeuro} (Causal VR Std.) & -- & $26.10 \pm 1.28$ & $44.20 \pm 2.50$ & $15.40 \pm 0.70$ & $18.10 \pm 0.92$ & $5.45 \pm 0.36$ & -- & $\mathbf{0.05}$ (0f buf) \\
 & Kinematic-Filter (SG, Non-Causal) & -- & $24.42 \pm 1.15$ & $41.80 \pm 2.34$ & $14.20 \pm 0.62$ & $16.85 \pm 0.84$ & $4.82 \pm 0.31$ & -- & $25.08$ (3f buf) \\
 & BRITS \cite{cao2018brits} (Bidirectional) & 0.64M / 0.29G & $15.10 \pm 0.78$ & $24.60 \pm 1.42$ & $9.35 \pm 0.46$ & $7.92 \pm 0.38$ & $2.64 \pm 0.19$ & $0.38$ & $250.84$ (Full seq) \\
 & SAITS \cite{du2023saits} (Non-Causal Attn.) & 0.72M / 0.31G & $14.22 \pm 0.65$ & $22.45 \pm 1.18$ & $8.74 \pm 0.39$ & $7.40 \pm 0.35$ & $2.38 \pm 0.17$ & $0.35$ & $250.81$ (Full seq) \\
 & ST-GCN \cite{yan2018spatial} & 0.45M / 0.19G & $14.95 \pm 0.71$ & $24.10 \pm 1.29$ & $9.10 \pm 0.42$ & $7.15 \pm 0.33$ & $2.72 \pm 0.21$ & $0.39$ & $100.85$ (Offline) \\
 & MS-G3D \cite{liu2020msg3d} & 0.82M / 0.44G & $13.48 \pm 0.58$ & $20.85 \pm 1.04$ & $8.12 \pm 0.35$ & $6.48 \pm 0.29$ & $2.55 \pm 0.18$ & $0.68$ & $101.14$ (Offline) \\
 & CTR-GCN \cite{chen2021ctrgcn} & 0.54M / 0.23G & $12.96 \pm 0.52$ & $19.74 \pm 0.96$ & $7.85 \pm 0.33$ & $6.12 \pm 0.27$ & $2.51 \pm 0.16$ & $0.51$ & $100.97$ (Offline) \\
 & BlockGCN \cite{zhou2024blockgcn} & 0.49M / 0.21G & $12.35 \pm 0.49$ & $18.60 \pm 0.88$ & $7.45 \pm 0.31$ & $5.80 \pm 0.25$ & $2.34 \pm 0.15$ & $0.48$ & $100.94$ (Offline) \\
 & SAITS + CTR-GCN (2-Stage) & 1.26M / 0.54G & $10.45 \pm 0.42$ & $14.85 \pm 0.68$ & $6.65 \pm 0.27$ & $5.62 \pm 0.24$ & $1.92 \pm 0.14$ & $0.86$ & $251.32$ (Offline) \\
 & CTR-GCN + Bridge $\mathcal{B}_{\text{temp}}$ & 0.55M / 0.24G & $9.82 \pm 0.38$ & $13.60 \pm 0.61$ & $6.34 \pm 0.25$ & $5.38 \pm 0.22$ & $1.76 \pm 0.13$ & $0.53$ & $100.99$ (12f buf) \\
 & \textbf{System B: KAR-STGCN (Causal Stream)} & \textbf{0.42M / 0.18G} & $9.18 \pm 0.35$ & $12.35 \pm 0.53$ & $6.08 \pm 0.23$ & $4.38 \pm 0.18$ & $1.61 \pm 0.12$ & $\mathbf{0.42}$ & $\mathbf{0.88}$ (\textbf{0f buf}) \\
 & \textbf{System A: KAR-STGCN (Offline Win.)} & \textbf{0.42M / 0.18G} & $\mathbf{8.64 \pm 0.31}$ & $\mathbf{11.42 \pm 0.48}$ & $\mathbf{5.82 \pm 0.21}$ & $\mathbf{4.12 \pm 0.16}$ & $\mathbf{1.48 \pm 0.11}$ & $\mathbf{0.42}$ & $100.88$ (12f buf) \\
\midrule
\multirow{6}{*}{\shortstack[l]{SHREC'17\\(Level 4:\\Compound)}} 
 & SAITS \cite{du2023saits} & 0.72M / 0.31G & $19.38 \pm 0.88$ & $29.40 \pm 1.52$ & $12.45 \pm 0.58$ & $10.82 \pm 0.51$ & $3.45 \pm 0.25$ & $\mathbf{0.36}$ & $250.82$ (Full seq) \\
 & CTR-GCN \cite{chen2021ctrgcn} & 0.54M / 0.23G & $17.80 \pm 0.74$ & $26.15 \pm 1.28$ & $11.28 \pm 0.48$ & $8.85 \pm 0.39$ & $3.40 \pm 0.22$ & $0.52$ & $100.98$ (Offline) \\
 & BlockGCN \cite{zhou2024blockgcn} & 0.49M / 0.21G & $16.95 \pm 0.68$ & $24.80 \pm 1.18$ & $10.75 \pm 0.45$ & $8.32 \pm 0.36$ & $3.18 \pm 0.20$ & $0.49$ & $100.95$ (Offline) \\
 & SAITS + CTR-GCN (2-Stage) & 1.26M / 0.54G & $13.90 \pm 0.55$ & $19.45 \pm 0.86$ & $9.12 \pm 0.36$ & $7.65 \pm 0.33$ & $2.45 \pm 0.17$ & $0.88$ & $251.34$ (Offline) \\
 & \textbf{System B: KAR-STGCN (Causal Stream)} & \textbf{0.42M / 0.18G} & $11.96 \pm 0.43$ & $16.45 \pm 0.62$ & $8.28 \pm 0.31$ & $5.72 \pm 0.24$ & $2.05 \pm 0.15$ & $0.42$ & $\mathbf{0.88}$ (\textbf{0f buf}) \\
 & \textbf{System A: KAR-STGCN (Offline Win.)} & \textbf{0.42M / 0.18G} & $\mathbf{11.38 \pm 0.39}$ & $\mathbf{15.64 \pm 0.58}$ & $\mathbf{7.94 \pm 0.28}$ & $\mathbf{5.46 \pm 0.22}$ & $\mathbf{1.92 \pm 0.14}$ & $0.44$ & $100.90$ (12f buf) \\
\midrule
\multirow{3}{*}{\shortstack[l]{FPHA\\(Level 1--3)}} 
 & Kinematic-Filter (L1 / L2 / L3) & -- & $4.48$ / $14.60$ / $31.25$ & -- / $28.40$ / $48.90$ & $3.10$ / $9.20$ / $18.40$ & $3.15$ / $9.80$ / $20.10$ & $1.10$ / $3.25$ / $5.90$ & -- & $25.08$ (3f buf) \\
 & CTR-GCN \cite{chen2021ctrgcn} (L1 / L2 / L3) & 0.54M / 0.23G & $2.68$ / $9.34$ / $17.82$ & -- / $18.10$ / $26.40$ & $1.75$ / $5.85$ / $10.90$ & $1.52$ / $5.40$ / $9.85$ & $0.64$ / $2.10$ / $3.45$ & $0.53$ & $100.99$ (Offline) \\
 & \textbf{KAR-STGCN (Sys. A, L1 / L2 / L3)} & \textbf{0.42M / 0.18G} & $\mathbf{2.46}$ / $\mathbf{6.54}$ / $\mathbf{11.95}$ & -- / $\mathbf{11.20}$ / $\mathbf{15.12}$ & $\mathbf{1.54}$ / $\mathbf{4.05}$ / $\mathbf{7.32}$ & $\mathbf{1.20}$ / $\mathbf{3.12}$ / $\mathbf{5.64}$ & $\mathbf{0.48}$ / $\mathbf{1.35}$ / $\mathbf{2.08}$ & $0.43$ & $100.89$ (12f buf) \\
\midrule
\multirow{6}{*}{\shortstack[l]{FPHA\\(Level 4:\\Compound)}} 
 & SAITS \cite{du2023saits} & 0.72M / 0.31G & $27.90 \pm 1.24$ & $38.60 \pm 1.85$ & $16.80 \pm 0.76$ & $16.45 \pm 0.78$ & $4.95 \pm 0.34$ & $\mathbf{0.37}$ & $250.83$ (Full seq) \\
 & CTR-GCN \cite{chen2021ctrgcn} & 0.54M / 0.23G & $25.65 \pm 1.05$ & $34.12 \pm 1.54$ & $15.18 \pm 0.64$ & $13.65 \pm 0.61$ & $4.72 \pm 0.30$ & $0.54$ & $101.00$ (Offline) \\
 & BlockGCN \cite{zhou2024blockgcn} & 0.49M / 0.21G & $24.40 \pm 0.98$ & $32.45 \pm 1.44$ & $14.50 \pm 0.60$ & $12.90 \pm 0.57$ & $4.45 \pm 0.28$ & $0.50$ & $100.96$ (Offline) \\
 & SAITS + CTR-GCN (2-Stage) & 1.26M / 0.54G & $20.15 \pm 0.82$ & $25.80 \pm 1.15$ & $12.10 \pm 0.49$ & $11.20 \pm 0.48$ & $3.42 \pm 0.23$ & $0.91$ & $251.37$ (Offline) \\
 & \textbf{System B: KAR-STGCN (Causal Stream)} & \textbf{0.42M / 0.18G} & $17.18 \pm 0.58$ & $20.95 \pm 0.76$ & $10.35 \pm 0.39$ & $8.24 \pm 0.36$ & $2.82 \pm 0.19$ & $0.43$ & $\mathbf{0.89}$ (\textbf{0f buf}) \\
 & \textbf{System A: KAR-STGCN (Offline Win.)} & \textbf{0.42M / 0.18G} & $\mathbf{16.42 \pm 0.52}$ & $\mathbf{19.88 \pm 0.71}$ & $\mathbf{9.86 \pm 0.36}$ & $\mathbf{7.85 \pm 0.34}$ & $\mathbf{2.64 \pm 0.18}$ & $0.44$ & $100.90$ (12f buf) \\
\bottomrule
\end{tabular}%
}
\end{table*}

\subsection{Gesture Recognition and Real Leap Motion (LMDHG) Transfer}
Table \ref{tab:classification_sota} reports subject-independent gesture recognition across \textbf{SHREC'17 (14 and 28 classes)} and \textbf{FPHA (45 classes)}. Under Level 4 compound dropouts, System A and System B achieve \textbf{91.84\%} and \textbf{91.12\%} on SHREC'17-14, \textbf{86.42\%} and \textbf{85.60\%} on SHREC'17-28, and \textbf{82.76\%} and \textbf{81.90\%} on FPHA-45, outperforming BlockGCN \cite{zhou2024blockgcn} and SAITS+CTR-GCN.

Table \ref{tab:sim_to_real_lmdhg} evaluates transfer onto real first-generation Leap Motion streams from \textbf{LMDHG} \cite{boulahia2017dynamic} ($N=21$ subjects, 13 classes), separating \textbf{Held-Out Contiguous Masking Restoration} ($\text{MPJPE}_{\text{miss}}$ and BLCE on verified segments) from \textbf{Natural Unmasked Recognition} (Macro-F1 on raw unedited sensor streams). Pre-training on our six-type dropout simulator (Protocol 2) yields a zero-shot held-out $\text{MPJPE}_{\text{miss}}$ of \textbf{12.84 $\pm$ 0.56 mm} and $82.35\%$ natural Macro-F1, while fine-tuning on LMDHG training subjects (Protocol 3) achieves \textbf{9.62 $\pm$ 0.41 mm} held-out $\text{MPJPE}_{\text{miss}}$ and \textbf{89.45 $\pm$ 1.38\%} natural Macro-F1.

\begin{table*}[!t]
\centering
\caption{Subject-independent gesture classification performance ($\text{Macro-F1}\% \uparrow$, $\text{Mean} \pm \text{SD}$) under our unified 5-fold cross-subject protocol on SHREC'17 (14-class and 28-class) and FPHA (45-class).}
\label{tab:classification_sota}
\resizebox{\textwidth}{!}{%
\begin{tabular}{lcccccc}
\toprule
\multirow{2}{*}{\textbf{Model Architecture / System}} & \multicolumn{2}{c}{\textbf{SHREC'17 (14 Classes)}} & \multicolumn{2}{c}{\textbf{SHREC'17 (28 Classes)}} & \multicolumn{2}{c}{\textbf{FPHA (45 Action Classes)}} \\
\cmidrule(lr){2-3} \cmidrule(lr){4-5} \cmidrule(lr){6-7}
 & \textbf{Level 1 (Clean)} & \textbf{Level 4 (Dropout)} & \textbf{Level 1 (Clean)} & \textbf{Level 4 (Dropout)} & \textbf{Level 1 (Clean)} & \textbf{Level 4 (Dropout)} \\
\midrule
BiLSTM-Baseline & $88.40 \pm 1.25$ & $64.20 \pm 2.45$ & $83.10 \pm 1.42$ & $56.80 \pm 2.80$ & $80.25 \pm 1.65$ & $52.40 \pm 3.10$ \\
DD-Net (Hand Kinematic) \cite{yang2019ddnet} & $93.80 \pm 0.88$ & $74.50 \pm 1.75$ & $90.10 \pm 0.95$ & $68.20 \pm 2.05$ & $87.90 \pm 1.15$ & $64.80 \pm 2.30$ \\
ST-GCN \cite{yan2018spatial} & $92.85 \pm 0.92$ & $73.15 \pm 1.85$ & $88.75 \pm 1.08$ & $66.40 \pm 2.10$ & $86.40 \pm 1.25$ & $63.10 \pm 2.40$ \\
2s-AGCN \cite{shi2019twostream} & $94.10 \pm 0.84$ & $75.80 \pm 1.72$ & $90.45 \pm 0.96$ & $69.10 \pm 1.95$ & $88.20 \pm 1.18$ & $65.75 \pm 2.20$ \\
MS-G3D \cite{liu2020msg3d} & $95.20 \pm 0.76$ & $77.90 \pm 1.54$ & $91.80 \pm 0.88$ & $71.40 \pm 1.78$ & $89.65 \pm 1.05$ & $67.80 \pm 1.95$ \\
CTR-GCN \cite{chen2021ctrgcn} & $95.65 \pm 0.72$ & $78.60 \pm 1.48$ & $92.35 \pm 0.82$ & $72.85 \pm 1.65$ & $90.10 \pm 0.98$ & $68.40 \pm 1.88$ \\
BlockGCN \cite{zhou2024blockgcn} & $95.90 \pm 0.68$ & $80.15 \pm 1.38$ & $92.70 \pm 0.78$ & $74.20 \pm 1.55$ & $90.65 \pm 0.92$ & $69.90 \pm 1.76$ \\
SAITS + CTR-GCN (2-Stage) & $95.70 \pm 0.70$ & $85.40 \pm 1.15$ & $92.40 \pm 0.80$ & $79.80 \pm 1.28$ & $90.20 \pm 0.95$ & $75.60 \pm 1.45$ \\
\midrule
\textbf{System B: KAR-STGCN (Causal Stream)} & $96.15 \pm 0.66$ & $91.12 \pm 0.89$ & $92.85 \pm 0.76$ & $85.60 \pm 1.02$ & $91.10 \pm 0.90$ & $81.90 \pm 1.18$ \\
\textbf{System A: KAR-STGCN (Offline Win.)} & $\mathbf{96.42 \pm 0.64}$ & $\mathbf{91.84 \pm 0.85}$ & $\mathbf{93.18 \pm 0.74}$ & $\mathbf{86.42 \pm 0.96}$ & $\mathbf{91.45 \pm 0.88}$ & $\mathbf{82.76 \pm 1.12}$ \\
\bottomrule
\end{tabular}%
}
\end{table*}

\begin{table*}[!b]
\centering
\caption{Cross-dataset transfer onto real first-generation Leap Motion streams from LMDHG \cite{boulahia2017dynamic} ($N=21$ subjects, 13 classes). Held-out restoration metrics ($\text{MPJPE}_{\text{miss}}$, BLCE) are evaluated on contiguous masked segments of verified real tracking frames, whereas Macro-F1 is evaluated on natural unmasked sensor streams.}
\label{tab:sim_to_real_lmdhg}
\resizebox{\textwidth}{!}{%
\begin{tabular}{llccccc}
\toprule
\textbf{Transfer Protocol (Train $\rightarrow$ Test)} & \textbf{Regime} & \textbf{SAITS Held-Out $\text{MPJPE}_{\text{miss}}$} $\downarrow$ & \textbf{CTR-GCN Held-Out $\text{MPJPE}_{\text{miss}}$} $\downarrow$ & \textbf{KAR-STGCN Held-Out $\text{MPJPE}_{\text{miss}}$} $\downarrow$ & \textbf{KAR-STGCN Held-Out BLCE} $\downarrow$ & \textbf{Natural Stream Macro-F1 (\%)} $\uparrow$ \\
\midrule
1. Uncalibrated Gaussian Sim $\rightarrow$ LMDHG & Zero-Shot & $19.85 \pm 0.94$ & $21.45 \pm 1.05$ & $\mathbf{17.42 \pm 0.82}$ & $8.94 \pm 0.45$ & $74.60 \pm 2.85$ \\
2. Six-Type Dropout Sim $\rightarrow$ LMDHG & Zero-Shot & $16.12 \pm 0.78$ & $17.90 \pm 0.86$ & $\mathbf{12.84 \pm 0.56}$ & $6.18 \pm 0.31$ & $82.35 \pm 1.94$ \\
3. Six-Type Sim + LMDHG $\rightarrow$ LMDHG & Fine-Tuned & $13.45 \pm 0.62$ & $14.28 \pm 0.65$ & $\mathbf{9.62 \pm 0.41}$ & $\mathbf{4.45 \pm 0.22}$ & $\mathbf{89.45 \pm 1.38}$ \\
\bottomrule
\end{tabular}%
}
\end{table*}

\subsection{Natural Dropout Restoration on Orthogonal Dual-LMC2 Hardware}
To evaluate natural optical dropouts with an independent physical reference (addressing the lack of external ground truth in single-sensor LMDHG), Table \ref{tab:dual_lmc2_gt} evaluates our framework on the \textbf{Orthogonal Dual-LMC2 Benchmark} \cite{gilmartin2024duallmc2} ($N=25$ subjects, $15$ dynamic gesture classes, $375$ sequences). In this setup, Primary Sensor 1 (horizontal desktop placement) experiences natural self-occlusion and oblique whole-hand dropouts, while Orthogonal Sensor 2 (calibrated into Sensor 1's coordinate frame via Kabsch $\text{SE}(3)$ rigid alignment over high-confidence frames) provides independent 3D ground truth during Sensor 1 tracking loss. On natural unmasked dropouts, raw Sensor 1 streams exhibit $\text{MPJPE}_{\text{miss}} = 46.50\text{ mm}$ and $34.20\%$ JAVR. \textbf{System B (Causal KAR-STGCN)} reduces natural $\text{MPJPE}_{\text{miss}}$ to \textbf{10.85 $\pm$ 0.49 mm} ($92.85\%$ Macro-F1 at $0.88\text{ ms}$ end-to-end latency), while \textbf{System A (Offline KAR-STGCN)} achieves \textbf{10.12 $\pm$ 0.45 mm} ($\text{BLCE}=3.54\text{ mm}$, $\text{JAVR}=1.05\%$, and \textbf{93.60\%} Macro-F1), outperforming BlockGCN \cite{zhou2024blockgcn} ($17.20\text{ mm}$) and SAITS+CTR-GCN ($13.90\text{ mm}$) with high statistical significance ($N=25, \Delta\text{MPJPE}_{\text{all}} = -4.65 \pm 1.62\text{ mm}, 95\%\text{ CI } [-5.32, -3.98]\text{ mm}, t(24) = -14.35, \text{Cohen's } d = 2.87, p_{\text{Holm}} < 0.0001$).

\begin{table*}[!t]
\centering
\caption{Natural optical dropout restoration and gesture recognition on the Orthogonal Dual-LMC2 hardware benchmark \cite{gilmartin2024duallmc2} ($N=25$ subjects, 15 classes), using Kabsch-aligned Orthogonal Sensor 2 as the independent ground-truth reference during Primary Sensor 1 natural dropouts.}
\label{tab:dual_lmc2_gt}
\resizebox{\textwidth}{!}{%
\begin{tabular}{llcccccc}
\toprule
\textbf{Algorithmic Paradigm} & \textbf{Method / System} & \textbf{Natural $\text{MPJPE}_{\text{all}}$ (mm)} $\downarrow$ & \textbf{Natural $\text{MPJPE}_{\text{miss}}$ (mm)} $\downarrow$ & \textbf{BLCE (mm)} $\downarrow$ & \textbf{Retargeting JAVR (\%)} $\downarrow$ & \textbf{Macro-F1 (\%)} $\uparrow$ & \textbf{End-to-End Lat. (ms)} $\downarrow$ \\
\midrule
Raw Primary Sensor 1 & Unrestored LMC2 Sensor 1 (Raw SDK) & $22.84 \pm 1.12$ & $46.50 \pm 2.45$ & $19.40 \pm 0.95$ & $34.20 \pm 2.30\%$ & $68.40 \pm 2.15\%$ & -- \\
VR Industry Filter & $1\text{€}$ Filter \cite{casiez2012oneeuro} (Causal VR Std.) & $19.65 \pm 0.96$ & $38.90 \pm 2.05$ & $15.80 \pm 0.82$ & $27.50 \pm 1.95\%$ & $72.15 \pm 1.90\%$ & $\mathbf{0.05}$ (0f buf) \\
Time-Series Imputer & SAITS \cite{du2023saits} (Non-Causal Attn.) & $12.45 \pm 0.62$ & $19.80 \pm 1.04$ & $8.15 \pm 0.41$ & $12.80 \pm 1.05\%$ & $84.30 \pm 1.28\%$ & $250.81$ (Full seq) \\
Skeleton Action GCN & CTR-GCN \cite{chen2021ctrgcn} (ICCV 2021) & $11.80 \pm 0.58$ & $18.45 \pm 0.92$ & $6.90 \pm 0.35$ & $10.95 \pm 0.88\%$ & $85.60 \pm 1.18\%$ & $100.97$ (Offline) \\
Skeleton Action GCN & BlockGCN \cite{zhou2024blockgcn} (CVPR 2024) & $11.15 \pm 0.54$ & $17.20 \pm 0.85$ & $6.35 \pm 0.32$ & $9.85 \pm 0.82\%$ & $86.95 \pm 1.12\%$ & $100.94$ (Offline) \\
Two-Stage Hybrid & SAITS + CTR-GCN (2-Stage) & $9.48 \pm 0.45$ & $13.90 \pm 0.68$ & $5.40 \pm 0.27$ & $6.45 \pm 0.61\%$ & $89.20 \pm 0.96\%$ & $251.32$ (Offline) \\
\midrule
\multirow{2}{*}{\textbf{Kinematic-Graph (Ours)}} 
 & \textbf{System B: KAR-STGCN (Causal Stream)} & $7.62 \pm 0.34$ & $10.85 \pm 0.49$ & $3.82 \pm 0.17$ & $1.64 \pm 0.28\%$ & $92.85 \pm 0.78\%$ & $\mathbf{0.88}$ (\textbf{0f buf}) \\
 & \textbf{System A: KAR-STGCN (Offline Win.)} & $\mathbf{7.15 \pm 0.31}$ & $\mathbf{10.12 \pm 0.45}$ & $\mathbf{3.54 \pm 0.15}$ & $\mathbf{1.05 \pm 0.21\%}$ & $\mathbf{93.60 \pm 0.72\%}$ & $100.88$ (12f buf) \\
\bottomrule
\end{tabular}%
}
\end{table*}

% ============================================================================
\section{Ablation and Robustness Analysis}
% ============================================================================
\subsection{Normalization Generalization across Multi-Scale Perturbations}
Table \ref{tab:norm_ablation} evaluates eight normalization strategies across five test-time scale factors ($\times 0.70, 0.85, 1.00, 1.15, 1.40$). While raw coordinates with scale augmentation (Row 4) improve over unaugmented wrist-centering ($78.40\%$ vs. $51.20\%$ mean cross-scale F1), they still suffer a $7.80\%$ accuracy drop at extreme scales ($\times 0.70$ and $\times 1.40$) and higher reconstruction error ($14.12\text{ mm}$) because network capacity is expended on learning scale invariance rather than fine finger articulation. Both \textbf{System B (Causal Running-Median)} and \textbf{System A (Offline Sequence-Median)} maintain consistent Macro-F1 across all five scales ($91.15\%$ and $91.62\%$) while eliminating division-by-$\epsilon$ explosions.

\subsection{Visibility Mask Corruption ($0\%\text{--}20\%$) and Extended Burst Sweep ($5\text{--}40$ Frames)}
Table \ref{tab:ablation_sensitivity} (Left) evaluates architectural ablations and five levels of visibility mask corruption: \textbf{Clean ($0\%$ error)}, \textbf{Low ($5\%$ flip error)}, \textbf{Medium ($10\%$)}, \textbf{High ($15\%$)}, and \textbf{Severe ($20\%$)}, alongside our \textbf{Inferred Kinematic Mask ($\hat{\mathbf{V}}$)} and \textbf{No Visibility Mask ($\mathbf{V}=\mathbf{1}$)}. Using the kinematic detector $\hat{\mathbf{V}}$ without ground-truth $\mathbf{V}$ increases $\text{MPJPE}_{\text{all}}$ by only $0.34\text{ mm}$ ($11.72\text{ mm}$ vs. $11.38\text{ mm}$), and under $15\%\text{--}20\%$ severe mask corruption, KAR-STGCN still outperforms SAITS+CTR-GCN with an oracle mask ($13.68\text{--}14.60\text{ mm}$ vs. $13.90\text{ mm}$).

Table \ref{tab:ablation_sensitivity} (Right) sweeps contiguous whole-hand burst lengths across \textbf{$5, 10, 15, 20, 25, 30, 35,$ and $40$ frames}, reporting $\text{MPJPE}_{\text{miss}}$, Macro-F1, and the \textbf{Restoration Failure Rate (\%)} (percentage of sequences where $\text{MPJPE}_{\text{miss}} > 25.0\text{ mm}$). Within the 25-frame receptive field ($5\text{--}25$ frames), KAR-STGCN keeps failure rates $\le 8.4\%$; beyond $30\text{--}40$ frames, error and failure rates rise sharply ($34.20\text{ mm}$, $68.5\%$ failure at 40 frames), transparently delineating the physical boundary of local temporal bridging.

\begin{table*}[!b]
\centering
\caption{Normalization generalization under SHREC'17 Level 4 across five explicit test-time hand scale factors ($\times 0.70, 0.85, 1.00, 1.15, 1.40$).}
\label{tab:norm_ablation}
\resizebox{\textwidth}{!}{%
\begin{tabular}{lccccccc}
\toprule
\textbf{Normalization Strategy} & \textbf{$\text{MPJPE}_{\text{all}}$ (mm)} $\downarrow$ & \textbf{BLCE (mm)} $\downarrow$ & \textbf{Scale $\times 0.70$ F1} & \textbf{Scale $\times 0.85$ F1} & \textbf{In-Scale $\times 1.00$ F1} & \textbf{Scale $\times 1.15$ F1} & \textbf{Scale $\times 1.40$ F1} \\
\midrule
1. Naive Instantaneous $L_2$ ($\epsilon = 10^{-6}$) & $8.78 \times 10^7$ & $1.24 \times 10^8$ & $56.10\%$ & $59.40\%$ & $61.40\%$ & $58.80\%$ & $55.30\%$ \\
2. Clamped Instantaneous $L_2$ ($\tau_{\min} = 15.0\text{ mm}$) & $18.92 \pm 0.84$ & $11.25 \pm 0.56$ & $80.80\%$ & $82.40\%$ & $83.45\%$ & $82.10\%$ & $80.75\%$ \\
3. Wrist-Centered Only (No Bone Scaling) & $14.85 \pm 0.62$ & $8.10 \pm 0.39$ & $38.40\%$ & $64.20\%$ & $84.10\%$ & $62.80\%$ & $36.50\%$ \\
4. Raw + Random Scale Augmentation ($\times 0.7\text{--}1.4$) & $14.12 \pm 0.58$ & $7.85 \pm 0.36$ & $74.20\%$ & $81.50\%$ & $86.20\%$ & $80.90\%$ & $73.10\%$ \\
5. Hand-Span Scaling ($\max_{i,j}\|\mathbf{p}_{t,i} - \mathbf{p}_{t,j}\|_2$) & $16.40 \pm 0.75$ & $9.42 \pm 0.48$ & $82.50\%$ & $84.10\%$ & $85.15\%$ & $83.90\%$ & $82.35\%$ \\
6. Unmasked Sequence-Median Palm Length & $13.65 \pm 0.51$ & $6.84 \pm 0.31$ & $87.40\%$ & $88.30\%$ & $88.90\%$ & $88.20\%$ & $87.60\%$ \\
\midrule
\textbf{7. System B: Causal Running-Median $L_2$ (Ours)} & $11.96 \pm 0.43$ & $5.72 \pm 0.24$ & $90.95\%$ & $91.20\%$ & $91.12\%$ & $91.25\%$ & $91.08\%$ \\
\textbf{8. System A: Offline Sequence-Median $L_2$ (Ours)} & $\mathbf{11.38 \pm 0.39}$ & $\mathbf{5.46 \pm 0.22}$ & $\mathbf{91.50\%}$ & $\mathbf{91.70\%}$ & $\mathbf{91.84\%}$ & $\mathbf{91.68\%}$ & $\mathbf{91.42\%}$ \\
\bottomrule
\end{tabular}%
}
\end{table*}

\begin{table*}[!t]
\centering
\caption{Ablation and robustness study on SHREC'17 Level 4: (Left) architectural ablations and 5-level visibility mask corruption ($0\%\text{--}20\%$); (Right) extended whole-hand burst length sweep ($5\text{--}40$ frames) with Restoration Failure Rate ($\%\text{ sequences with MPJPE}_{\text{miss}} > 25\text{ mm}$).}
\label{tab:ablation_sensitivity}
\resizebox{\textwidth}{!}{%
\begin{tabular}{lccc|lcccc}
\toprule
\textbf{Architecture / Visibility Mask Regime} & \textbf{$\text{MPJPE}_{\text{all}}$} $\downarrow$ & \textbf{BLCE} $\downarrow$ & \textbf{Macro-F1 (\%)} $\uparrow$ & \textbf{Burst Length (Frames)} & \textbf{SAITS+CTR-GCN $\text{MPJPE}_{\text{miss}}$} & \textbf{KAR-STGCN $\text{MPJPE}_{\text{miss}}$} & \textbf{KAR-STGCN F1 (\%)} & \textbf{Failure Rate (\%)} $\downarrow$ \\
\midrule
w/o Dilated Bridge $\mathcal{B}_{\text{temp}}$ & $18.45 \pm 0.78$ & $9.12 \pm 0.42$ & $83.20 \pm 1.35$ & 5 Frames ($42\text{ ms}$) & $10.20 \pm 0.45$ & $\mathbf{7.85 \pm 0.32}$ & $\mathbf{94.80\%}$ & $\mathbf{0.0\%}$ \\
Narrow Bridge ($K_b=5, d=1$, RF=5) & $14.92 \pm 0.56$ & $6.95 \pm 0.31$ & $88.10 \pm 1.08$ & 10 Frames ($83\text{ ms}$) & $14.10 \pm 0.62$ & $\mathbf{10.95 \pm 0.44}$ & $\mathbf{93.25\%}$ & $\mathbf{0.4\%}$ \\
w/o Bone Loss ($\lambda_{\text{bone}}=0$) & $12.65 \pm 0.46$ & $10.54 \pm 0.49$ & $90.15 \pm 0.94$ & 15 Frames ($125\text{ ms}$) & $18.90 \pm 0.85$ & $\mathbf{14.20 \pm 0.56}$ & $\mathbf{91.10\%}$ & $\mathbf{1.8\%}$ \\
w/o Velocity Stream ($\Delta\hat{\mathbf{X}}_t$) & $11.52 \pm 0.41$ & $5.54 \pm 0.23$ & $86.40 \pm 1.18$ & 20 Frames ($167\text{ ms}$) & $23.15 \pm 1.08$ & $\mathbf{17.30 \pm 0.71}$ & $\mathbf{88.75\%}$ & $\mathbf{4.2\%}$ \\
No Visibility Mask ($\mathbf{V}=\mathbf{1}$ Blind) & $15.80 \pm 0.64$ & $7.45 \pm 0.35$ & $85.90 \pm 1.22$ & 25 Frames ($208\text{ ms}$, RF) & $27.40 \pm 1.30$ & $\mathbf{20.45 \pm 0.88}$ & $\mathbf{85.60\%}$ & $\mathbf{8.4\%}$ \\
Inferred Kinematic Mask $\hat{\mathbf{V}}$ (No GT) & $11.72 \pm 0.43$ & $5.62 \pm 0.24$ & $91.40 \pm 0.89$ & 30 Frames ($250\text{ ms}$, $>$RF) & $32.10 \pm 1.52$ & $\mathbf{24.60 \pm 1.12}$ & $\mathbf{81.20\%}$ & $24.5\%$ \\
Mask Error: Low ($5\%$) / Med ($10\%$) & $12.10$ / $12.85$ & $5.82$ / $6.15$ & $90.85\%$ / $89.70\%$ & 35 Frames ($292\text{ ms}$, $>$RF) & $36.80 \pm 1.75$ & $\mathbf{28.90 \pm 1.35}$ & $\mathbf{76.40\%}$ & $46.2\%$ \\
Mask Error: High ($15\%$) / Sev ($20\%$) & $13.68$ / $14.60$ & $6.58$ / $7.10$ & $88.30\%$ / $86.80\%$ & 40 Frames ($333\text{ ms}$, $>$RF) & $41.50 \pm 2.05$ & $\mathbf{34.20 \pm 1.64}$ & $\mathbf{70.15\%}$ & $68.5\%$ \\
\bottomrule
\end{tabular}%
}
\end{table*}

\subsection{True Subject-Level Statistical Hypothesis Testing}
To eliminate seed-level pseudoreplication (particularly on FPHA where $N=6$ subjects), Table \ref{tab:anova_results} reports statistical tests strictly using \textbf{true independent subjects} as the unit of analysis ($N=28, df=27$ on SHREC'17; $N=6, df=5$ on FPHA by averaging across the 5 random seeds per subject prior to testing; $N=21, df=20$ on LMDHG; and $N=25, df=24$ on Dual-LMC2). For FPHA, we additionally verify significance using a Linear Mixed-Effects Model ($\text{LMM } F_{1,23}$) with random intercepts for subject and seed. On FPHA Level 4 ($N=6, df=5$), KAR-STGCN achieves a mean error reduction of $\Delta\text{MPJPE} = -9.23\text{ mm}$ ($\text{SD}_{\text{diff}}=2.33\text{ mm}, \text{SE}=0.95\text{ mm}$), yielding paired $t(5) = -9.71$, two-sided $p = 0.0002$ ($p_{\text{Holm}} = 0.0012$, exact Wilcoxon signed-rank $p = 0.0312$, and LMM $p < 10^{-6}$), confirming statistically valid subject-level generalization.

\begin{table*}[!b]
\centering
\caption{True subject-level statistical hypothesis testing on $\text{MPJPE}_{\text{all}}$ (mm): Omnibus Welch's ANOVA ($F_W$), effect size ($\eta^2$), and Holm-Bonferroni corrected paired comparison of KAR-STGCN vs. CTR-GCN \cite{chen2021ctrgcn} using exact independent subject counts ($N=28$ on SHREC'17, $N=6$ on FPHA, $N=21$ on LMDHG, and $N=25$ on Dual-LMC2).}
\label{tab:anova_results}
\resizebox{\textwidth}{!}{%
\begin{tabular}{llccccccc}
\toprule
\textbf{Dataset (Independent Subjects $N, df$)} & \textbf{Degradation / Transfer Level} & \textbf{Welch's $F_W$} & \textbf{Effect $\eta^2$} & \textbf{Mean $\Delta\text{MPJPE} \pm \text{SD}_{\text{diff}}$} & \textbf{95\% CI of Difference} & \textbf{Paired $t(df)$} & \textbf{Cohen's $d$} & \textbf{$p_{\text{Holm}}$ (Wilcoxon $p$)} \\
\midrule
\multirow{2}{*}{SHREC'17 ($N=28$ subjects, $df=27$)} 
 & Level 3: Whole-Hand Burst & $142.65$ & $0.86$ & $-4.32 \pm 1.74\text{ mm}$ & $[-5.00, -3.64]\text{ mm}$ & $-13.12$ & $2.48$ & $<0.0001^*$ ($<0.0001$) \\
 & Level 4: Compound Long-Tail & $168.40$ & $0.88$ & $-6.42 \pm 2.26\text{ mm}$ & $[-7.30, -5.54]\text{ mm}$ & $-15.03$ & $2.84$ & $<0.0001^*$ ($<0.0001$) \\
\midrule
\multirow{2}{*}{FPHA ($N=6$ subjects, $df=5$)$^\dagger$} 
 & Level 3: Whole-Hand Burst & $48.20$ & $0.84$ & $-5.87 \pm 1.81\text{ mm}$ & $[-7.77, -3.97]\text{ mm}$ & $-7.94$ & $3.24$ & $0.0020^*$ ($0.0312$) \\
 & Level 4: Compound Long-Tail & $62.50$ & $0.87$ & $-9.23 \pm 2.33\text{ mm}$ & $[-11.67, -6.79]\text{ mm}$ & $-9.71$ & $3.97$ & $0.0012^*$ ($0.0312$) \\
\midrule
\multirow{2}{*}{LMDHG ($N=21$ subjects, $df=20$)} 
 & Protocol 2: Sim $\rightarrow$ Real (Zero-Shot) & $134.80$ & $0.85$ & $-5.06 \pm 2.14\text{ mm}$ & $[-6.04, -4.08]\text{ mm}$ & $-10.81$ & $2.36$ & $<0.0001^*$ ($<0.0001$) \\
 & Protocol 3: Sim $\rightarrow$ Real (Fine-Tuned) & $162.10$ & $0.88$ & $-4.66 \pm 1.72\text{ mm}$ & $[-5.44, -3.88]\text{ mm}$ & $-12.42$ & $2.71$ & $<0.0001^*$ ($<0.0001$) \\
\midrule
Dual-LMC2 ($N=25$ subjects, $df=24$) & Natural Cross-View GT Restoration & $156.40$ & $0.87$ & $-4.65 \pm 1.62\text{ mm}$ & $[-5.32, -3.98]\text{ mm}$ & $-14.35$ & $2.87$ & $<0.0001^*$ ($<0.0001$) \\
\bottomrule
\multicolumn{9}{l}{\footnotesize $*$ Statistically significant after Holm-Bonferroni correction. $^\dagger$On FPHA ($N=6$), exact two-sided Wilcoxon signed-rank minimum achievable $p=2/2^6=0.0312$; hierarchical LMM across $6\times 5$ runs yields $p<10^{-6}$.}
\end{tabular}%
}
\end{table*}

\begin{figure*}[!t]
\centering
\setlength{\unitlength}{1mm}
\begin{picture}(174, 52)
  % Panel (a): Ground-Truth 21-Joint Pose
  \put(2, 47){\scriptsize\textbf{(a) Ground-Truth Skeleton}}
  \put(1, 5){\framebox(40, 40){}}
  \put(21, 9){\textcolor{cGreen}{\circle*{2.0}}}
  \put(21, 9){\textcolor{cGreen}{\line(-1,1){10}}}
  \put(11, 19){\textcolor{cGreen}{\circle*{1.6}}}
  \put(11, 19){\textcolor{cGreen}{\line(-1,2){4}}}
  \put(7, 27){\textcolor{cGreen}{\circle*{1.6}}}
  \put(21, 9){\textcolor{cGreen}{\line(-1,3){6}}}
  \put(15, 27){\textcolor{cGreen}{\circle*{1.6}}}
  \put(15, 27){\textcolor{cGreen}{\line(0,1){11}}}
  \put(15, 38){\textcolor{cGreen}{\circle*{1.6}}}
  \put(21, 9){\textcolor{cGreen}{\line(0,1){20}}}
  \put(21, 29){\textcolor{cGreen}{\circle*{1.8}}}
  \put(21, 29){\textcolor{cGreen}{\line(0,1){11}}}
  \put(21, 40){\textcolor{cGreen}{\circle*{1.6}}}
  \put(21, 9){\textcolor{cGreen}{\line(1,3){6}}}
  \put(27, 27){\textcolor{cGreen}{\circle*{1.6}}}
  \put(27, 27){\textcolor{cGreen}{\line(0,1){11}}}
  \put(27, 38){\textcolor{cGreen}{\circle*{1.6}}}
  \put(21, 9){\textcolor{cGreen}{\line(1,2){9}}}
  \put(30, 27){\textcolor{cGreen}{\circle*{1.6}}}
  \put(30, 27){\textcolor{cGreen}{\line(1,2){4}}}
  \put(34, 35){\textcolor{cGreen}{\circle*{1.6}}}
  \put(3, 1){\scriptsize Palm Ref. $\|\mathbf{p}_9-\mathbf{p}_0\|_2$}

  % Panel (b): Corrupted Input (Dropout & Drift)
  \put(46, 47){\scriptsize\textbf{(b) Corrupted Input ($24.8\text{ mm}$)}}
  \put(45, 5){\framebox(40, 40){}}
  \put(65, 9){\textcolor{cGray}{\circle*{2.0}}}
  \put(65, 9){\textcolor{cGray}{\line(-1,1){10}}}
  \put(55, 19){\textcolor{cRed}{\circle*{1.8}}}
  \put(49, 15){\textcolor{cRed}{\textbf{$\times$}}}
  \put(65, 9){\textcolor{cGray}{\line(-1,3){6}}}
  \put(59, 27){\textcolor{cGray}{\circle*{1.6}}}
  \put(59, 27){\textcolor{cRed}{\line(-1,-2){4}}}
  \put(55, 19){\textcolor{cRed}{\textbf{$\times$}}}
  \put(65, 9){\textcolor{cGray}{\line(0,1){20}}}
  \put(65, 29){\textcolor{cGray}{\circle*{1.8}}}
  \put(65, 29){\textcolor{cRed}{\line(1,-2){4}}}
  \put(69, 21){\textcolor{cRed}{\textbf{$\times$}}}
  \put(65, 9){\textcolor{cGray}{\line(1,3){6}}}
  \put(71, 27){\textcolor{cGray}{\circle*{1.6}}}
  \put(71, 27){\textcolor{cRed}{\line(1,-1){5}}}
  \put(76, 22){\textcolor{cRed}{\textbf{$\times$}}}
  \put(65, 9){\textcolor{cGray}{\line(1,2){9}}}
  \put(74, 27){\textcolor{cRed}{\textbf{$\times$}}}
  \put(47, 1){\scriptsize \textcolor{cRed}{$\times$ Dropped / Drifted Joints}}

  % Panel (c): Restored KAR-STGCN with Heatmap
  \put(90, 47){\scriptsize\textbf{(c) KAR-STGCN ($6.8\text{ mm}$)}}
  \put(89, 5){\framebox(40, 40){}}
  \put(109, 9){\textcolor{cBlue}{\circle*{2.0}}}
  \put(109, 9){\textcolor{cBlue}{\line(-1,1){10}}}
  \put(99, 19){\textcolor{cBlue}{\circle*{1.6}}}
  \put(99, 19){\textcolor{cBlue}{\line(-1,2){4}}}
  \put(95, 27){\textcolor{cOrange}{\circle*{1.9}}}
  \put(109, 9){\textcolor{cBlue}{\line(-1,3){6}}}
  \put(103, 27){\textcolor{cBlue}{\circle*{1.6}}}
  \put(103, 27){\textcolor{cBlue}{\line(0,1){11}}}
  \put(103, 38){\textcolor{cOrange}{\circle*{1.9}}}
  \put(109, 9){\textcolor{cBlue}{\line(0,1){20}}}
  \put(109, 29){\textcolor{cBlue}{\circle*{1.8}}}
  \put(109, 29){\textcolor{cBlue}{\line(0,1){11}}}
  \put(109, 40){\textcolor{cOrange}{\circle*{1.9}}}
  \put(109, 9){\textcolor{cBlue}{\line(1,3){6}}}
  \put(115, 27){\textcolor{cBlue}{\circle*{1.6}}}
  \put(115, 27){\textcolor{cBlue}{\line(0,1){11}}}
  \put(115, 38){\textcolor{cOrange}{\circle*{1.9}}}
  \put(109, 9){\textcolor{cBlue}{\line(1,2){9}}}
  \put(118, 27){\textcolor{cBlue}{\circle*{1.6}}}
  \put(118, 27){\textcolor{cBlue}{\line(1,2){4}}}
  \put(122, 35){\textcolor{cOrange}{\circle*{1.9}}}
  \put(90, 1){\scriptsize \textcolor{cBlue}{Heatmap: Blue $<4$, Org $\approx 8\text{ mm}$}}

  % Panel (d): Failure Case (>35f Burst)
  \put(134, 47){\scriptsize\textbf{(d) Failure Case (38f Burst)}}
  \put(133, 5){\framebox(40, 40){}}
  \put(153, 9){\textcolor{cPurple}{\circle*{2.0}}}
  \put(153, 9){\textcolor{cPurple}{\line(-1,1){8}}}
  \put(145, 17){\textcolor{cPurple}{\circle*{1.6}}}
  \put(145, 17){\textcolor{cPurple}{\line(-1,1){3}}}
  \put(142, 20){\textcolor{cRed}{\circle*{1.9}}}
  \put(153, 9){\textcolor{cPurple}{\line(-1,3){4}}}
  \put(149, 21){\textcolor{cPurple}{\circle*{1.6}}}
  \put(149, 21){\textcolor{cPurple}{\line(0,1){6}}}
  \put(149, 27){\textcolor{cRed}{\circle*{1.9}}}
  \put(153, 9){\textcolor{cPurple}{\line(0,1){13}}}
  \put(153, 22){\textcolor{cPurple}{\circle*{1.8}}}
  \put(153, 22){\textcolor{cPurple}{\line(0,1){6}}}
  \put(153, 28){\textcolor{cRed}{\circle*{1.9}}}
  \put(153, 9){\textcolor{cPurple}{\line(1,3){4}}}
  \put(157, 21){\textcolor{cPurple}{\circle*{1.6}}}
  \put(157, 21){\textcolor{cPurple}{\line(0,1){6}}}
  \put(157, 27){\textcolor{cRed}{\circle*{1.9}}}
  \put(153, 9){\textcolor{cPurple}{\line(1,2){6}}}
  \put(159, 21){\textcolor{cPurple}{\circle*{1.6}}}
  \put(159, 21){\textcolor{cPurple}{\line(1,1){3}}}
  \put(162, 24){\textcolor{cRed}{\circle*{1.9}}}
  \put(134, 1){\scriptsize \textcolor{cRed}{Mean-Pose Collapse ($>25\text{f RF}$)}}
\end{picture}
\caption{Qualitative 3D hand restoration analysis: (a) Ground-truth 21-joint pose, (b) Corrupted observation (12-frame burst dropout), (c) KAR-STGCN restored skeleton with per-joint error heatmap (blue: $<4\text{ mm}$ on palm/MCP, orange: $6\text{--}9\text{ mm}$ on distal tips), and (d) Failure case under an extreme 38-frame burst exceeding the 25-frame bridge receptive field, resulting in regression toward a semi-flexed mean hand pose.}
\label{fig:qualitative_hand}
\end{figure*}

% ============================================================================
\section{Real-World Multimedia Application: KAR-Bench Studio}
% ============================================================================
To demonstrate direct utility for multimedia production and VR/AR systems, we developed \textbf{KAR-Bench Studio}, an open-source diagnostic and 3D character retargeting suite. KAR-Bench Studio integrates: (1) a live palm-bone oscilloscope monitoring $\|\mathbf{p}_{t,9}-\mathbf{p}_{t,0}\|_2$ stability, (2) a TorchScript deployment engine ($1.72\text{ MB}$ serialized model), and (3) an automated Biovision Hierarchy (\texttt{.bvh}) skeletal exporter compatible with Blender, Unity, and Unreal Engine.

Table \ref{tab:finger_retargeting} decomposes restoration errors across the five fingers on SHREC'17 Level 4 and reports the downstream \textbf{Joint-Angle Violation Rate (JAVR \%)}—the percentage of frames where any interphalangeal joint violates human biomechanical limits (hyperextension $<-15^\circ$ or lateral abduction $>25^\circ$) during 3D character rig retargeting. Because the ring and pinky fingers suffer the highest self-occlusion rates, the $1\text{€}$ Filter \cite{casiez2012oneeuro} and RTS Kalman Smoother exhibit severe pinky errors ($28.65\text{ mm}$ and $22.30\text{ mm}$) and high JAVR ($31.40\%$ and $19.85\%$). Both \textbf{System B (Causal KAR-STGCN)} and \textbf{System A (Offline KAR-STGCN)} leverage cross-finger graph adjacency ($\hat{\mathbf{A}}$) and $\mathcal{L}_{\text{bone}}$ to constrain pinky error to $11.28\text{ mm}$ and $10.42\text{ mm}$ while reducing JAVR to \textbf{1.85\%} and \textbf{1.12\%}, respectively.

\begin{table*}[!t]
\centering
\caption{Cross-paradigm anatomical finger-wise error breakdown ($\text{MPJPE}$ in mm $\downarrow$), downstream 3D character retargeting Joint-Angle Violation Rate ($\text{JAVR}\% \downarrow$), and CPU/GPU forward latency on SHREC'17 Level 4.}
\label{tab:finger_retargeting}
\resizebox{\textwidth}{!}{%
\begin{tabular}{llcccccccc}
\toprule
\textbf{Algorithmic Paradigm} & \textbf{Method / System} & \textbf{Thumb (J1--4)} & \textbf{Index (J5--8)} & \textbf{Middle (J9--12)} & \textbf{Ring (J13--16)} & \textbf{Pinky (J17--20)} & \textbf{Retargeting JAVR (\%)} $\downarrow$ & \textbf{CPU Lat. (ms)} & \textbf{GPU Fwd (ms)} \\
\midrule
VR Industry Filter & $1\text{€}$ Filter \cite{casiez2012oneeuro} & $18.42 \pm 0.85$ & $22.15 \pm 1.04$ & $21.80 \pm 0.98$ & $25.40 \pm 1.22$ & $28.65 \pm 1.45$ & $31.40 \pm 2.15\%$ & $\mathbf{0.05}$ & -- \\
Kinematic Smoother & RTS Kalman Smoother & $14.20 \pm 0.68$ & $17.45 \pm 0.82$ & $16.90 \pm 0.79$ & $19.85 \pm 0.94$ & $22.30 \pm 1.12$ & $19.85 \pm 1.64\%$ & $0.18$ & -- \\
Time-Series Imputer & SAITS \cite{du2023saits} & $11.85 \pm 0.54$ & $14.10 \pm 0.65$ & $13.65 \pm 0.62$ & $16.20 \pm 0.78$ & $18.45 \pm 0.89$ & $14.20 \pm 1.18\%$ & $3.42$ & $\mathbf{0.36}$ \\
Skeleton Action GCN & BlockGCN \cite{zhou2024blockgcn} & $10.90 \pm 0.49$ & $13.25 \pm 0.61$ & $12.80 \pm 0.58$ & $15.10 \pm 0.72$ & $17.20 \pm 0.84$ & $11.65 \pm 0.95\%$ & $4.85$ & $0.49$ \\
Two-Stage Hybrid & SAITS + CTR-GCN & $9.45 \pm 0.42$ & $11.20 \pm 0.51$ & $10.85 \pm 0.48$ & $12.60 \pm 0.59$ & $14.15 \pm 0.68$ & $7.80 \pm 0.72\%$ & $8.10$ & $0.88$ \\
\midrule
\multirow{2}{*}{\textbf{Kinematic-Graph (Ours)}} 
 & \textbf{System B: KAR-STGCN (Causal)} & $7.85 \pm 0.34$ & $9.12 \pm 0.41$ & $8.75 \pm 0.38$ & $10.15 \pm 0.46$ & $11.28 \pm 0.52$ & $1.85 \pm 0.31\%$ & $1.82$ & $0.42$ \\
 & \textbf{System A: KAR-STGCN (Offline)} & $\mathbf{7.24 \pm 0.31}$ & $\mathbf{8.45 \pm 0.38}$ & $\mathbf{8.10 \pm 0.35}$ & $\mathbf{9.38 \pm 0.42}$ & $\mathbf{10.42 \pm 0.48}$ & $\mathbf{1.12 \pm 0.24\%}$ & $1.85$ & $0.44$ \\
\bottomrule
\end{tabular}%
}
\end{table*}

% ============================================================================
\section{Discussion, Limitations, and Failure Cases}
% ============================================================================
\begin{enumerate}
    \item \textbf{Failure Mode Analysis under Prolonged Dropouts ($>25$ Frames):} As illustrated in Figure \ref{fig:qualitative_hand}(d) and quantified in Table \ref{tab:ablation_sensitivity} (Right), when a whole-hand dropout exceeds the 25-frame receptive field of $\mathcal{B}_{\text{temp}}$ (e.g., $30\text{--}40$ frames), the temporal bridge loses boundary anchor frames and regresses toward a static, semi-flexed mean hand posture during the middle of the gap. Recovering from prolonged hand exits ($>0.3\text{ s}$) requires generative motion priors \cite{rastgoo2025deep} rather than local convolutional bridging.
    \item \textbf{Cross-Sensor Ground Truth vs. Multi-Camera OptiTrack Studio:} By incorporating the Orthogonal Dual-LMC2 benchmark \cite{gilmartin2024duallmc2} (Table \ref{tab:dual_lmc2_gt}), we validate restoration on natural LMC2 dropouts using a Kabsch-aligned secondary optical sensor. Nevertheless, stereo-IR cross-view ground truth carries an intrinsic $\approx 1\text{--}2\text{ mm}$ baseline variance compared to studio multi-camera marker arrays.
    \item \textbf{Reference Modality Differences:} SHREC'17 skeletons are derived from depth images \cite{de2017shrec} and thus represent an approximate reference, whereas FPHA \cite{garcia2018first} provides direct 6-DoF magnetic sensor ground truth.
\end{enumerate}

% ============================================================================
\section{Conclusion}
% ============================================================================
This paper introduced KAR-STGCN, a visibility-guarded $L_2$ palm-bone normalization scheme (in both offline non-causal and causal streaming regimes), and the KAR-Bench Studio retargeting suite for 3D hand skeletal restoration and gesture recognition. By preventing division-by-zero singularities, mitigating scale shortcuts, and enforcing anatomical bone consistency across 25-frame dilated neighborhoods, KAR-STGCN significantly outperforms VR filters, time-series imputers, and skeleton GCNs across SHREC'17, FPHA, LMDHG, and Dual-LMC2 benchmarks.

% ============================================================================
\section*{Declarations}
% ============================================================================
\begin{itemize}
    \item \textbf{Funding:} This research project was financially supported by the Faculty of Informatics, Mahasarakham University.
    \item \textbf{Conflict of Interest:} The author declares no conflict of interest.
    \item \textbf{Author's Contributions:} The sole author conceived the study, designed the methodology, implemented the software pipeline, conducted the experiments, performed the statistical analysis, and wrote the manuscript.
    \item \textbf{Use of AI Tools:} Generative AI tools were utilized strictly for English language polishing and LaTeX table formatting verification, in accordance with Springer Nature authorship guidelines. All experimental code, numerical data, and scientific conclusions were verified by the author.
    \item \textbf{Data and Code Availability:} The source code, KAR-Bench Studio software suite, multi-seed evaluation scripts, and benchmark protocols are publicly available at \url{https://github.com/siripornp-msu/KAR-STGCN-3D-Hand-Restoration}.
\end{itemize}

\begin{thebibliography}{99}

\bibitem{bregler2007motion}
C. Bregler, ``Motion capture technology for entertainment [In the Spotlight],'' \textit{IEEE Signal Processing Magazine}, vol. 24, no. 6, pp. 160--158, 2007.

\bibitem{tirakoat2011optimized}
S. Tirakoat, ``Optimized motion capture system for full body human motion capturing: Case study of educational institution and small animation production,'' in \textit{2011 Workshop on Digital Media and Digital Content Management (DMDCM)}, IEEE, 2011, pp. 117--120, doi: 10.1109/DMDCM.2011.33.

\bibitem{noiumkar2013use}
S. Noiumkar and S. Tirakoat, ``Use of optical motion capture in sports science: A case study of golf swing,'' in \textit{2013 International Conference on Informatics and Creative Multimedia (ICICM)}, IEEE, 2013, pp. 310--313, doi: 10.1109/ICICM.2013.58.

\bibitem{marin2016hand}
G. Marin, F. Dominio, and P. Zanuttigh, ``Hand gesture recognition with jointly calibrated Leap Motion and depth sensor,'' \textit{Multimedia Tools and Applications}, vol. 75, no. 22, pp. 14991--15015, 2016, doi: 10.1007/s11042-015-2451-6.

\bibitem{boulahia2017dynamic}
S. Y. Boulahia, E. Anquetil, F. Multon, and R. Kulpa, ``Dynamic hand gesture recognition based on 3D pattern assembled trajectories,'' in \textit{7th International Conference on Image Processing Theory, Tools and Applications (IPTA)}, IEEE, 2017, pp. 1--6.

\bibitem{de2017shrec}
Q. De Smedt, H. Wannous, J.-P. Vandeborre, J. Guerry, B. Le Saux, and D. Filliat, ``SHREC'17 track: 3D hand gesture recognition using a depth and skeletal dataset,'' in \textit{Eurographics Workshop on 3D Object Retrieval (3DOR)}, 2017, pp. 1--6.

\bibitem{garcia2018first}
G. Garcia-Hernando, S. Yuan, S. Baek, and T.-K. Kim, ``First-person hand action benchmark with RGB-D videos and 3D hand pose annotations,'' in \textit{Proceedings of the IEEE Conference on Computer Vision and Pattern Recognition (CVPR)}, 2018, pp. 409--419.

\bibitem{kumar2018position}
P. Kumar, R. Saini, P. P. Roy, and D. P. Dogra, ``A position and rotation invariant framework for sign language recognition (SLR) using Kinect,'' \textit{Multimedia Tools and Applications}, vol. 77, no. 7, pp. 8823--8846, 2018, doi: 10.1007/s11042-017-4776-9.

\bibitem{yan2018spatial}
S. Yan, Y. Xiong, and D. Lin, ``Spatial temporal graph convolutional networks for skeleton-based action recognition,'' in \textit{Proceedings of the AAAI Conference on Artificial Intelligence}, vol. 32, no. 1, 2018.

\bibitem{cao2018brits}
W. Cao, D. Wang, J. Li, H. Zhou, L. Li, and Y. Li, ``BRITS: Bidirectional recurrent imputation for time series,'' in \textit{Advances in Neural Information Processing Systems (NeurIPS)}, vol. 31, 2018.

\bibitem{shi2019twostream}
L. Shi, Y. Zhang, J. Cheng, and H. Lu, ``Two-stream adaptive graph convolutional networks for skeleton-based action recognition,'' in \textit{Proceedings of the IEEE/CVF Conference on Computer Vision and Pattern Recognition (CVPR)}, 2019, pp. 12026--12035.

\bibitem{yang2019ddnet}
F. Yang, Y. Wu, S. Sakti, and S. Nakamura, ``Make skeleton-based action recognition model smaller, faster and better,'' in \textit{Proceedings of the ACM Multimedia Asia}, 2019, pp. 1--6.

\bibitem{liu2020msg3d}
Z. Liu, H. Zhang, Z. Chen, Z. Wang, and W. Ouyang, ``Disentangling and unifying graph convolutions for skeleton-based action recognition,'' in \textit{Proceedings of the IEEE/CVF Conference on Computer Vision and Pattern Recognition (CVPR)}, 2020, pp. 143--152.

\bibitem{chen2021ctrgcn}
Y. Chen, Z. Zhang, C. Yuan, B. Li, Y. Deng, and W. Hu, ``Channel-wise topology refinement graph convolution for skeleton-based action recognition,'' in \textit{Proceedings of the IEEE/CVF International Conference on Computer Vision (ICCV)}, 2021, pp. 13359--13368.

\bibitem{asteasuain2022applying}
F. Asteasuain, ``Applying machine learning to Leap Motion hand gesture recognition,'' \textit{Multimedia Tools and Applications}, vol. 81, pp. 34211--34232, 2022.

\bibitem{du2023saits}
W. Du, D. C{\^o}t{\'e}, and Y. Liu, ``SAITS: Self-attention-based imputation for time series,'' \textit{Expert Systems with Applications}, vol. 219, p. 119619, 2023.

\bibitem{zhou2024blockgcn}
Y. Zhou, X. Yan, Z.-Q. Cheng, Y. Yan, Q. Dai, and X.-S. Hua, ``BlockGCN: Redefine topology awareness for skeleton-based action recognition,'' in \textit{Proceedings of the IEEE/CVF Conference on Computer Vision and Pattern Recognition (CVPR)}, 2024, pp. 2049--2058.

\bibitem{gilmartin2024duallmc2}
M. Gil-Mart{\'\i}n, R. San-Segundo, F. Fern{\'a}ndez-Mart{\'\i}nez, and J. Ferreiros-L{\'o}pez, ``Dual Leap Motion Controller 2 dataset for hand motion and gesture recognition,'' \textit{Scientific Data}, vol. 11, p. 1012, 2024, doi: 10.1038/s41597-024-03968-9.

\bibitem{liu2024novel}
Y. Liu, X. Li, H. Yu, and L. Yang, ``A novel sEMG-based dynamic hand gesture recognition approach via residual attention network,'' \textit{Multimedia Tools and Applications}, vol. 83, no. 4, pp. 9329--9349, 2024, doi: 10.1007/s11042-023-15748-5.

\bibitem{matuszczyk2025lmc2}
F. Matuszczyk, M. J{\k{e}}drusiak, J. Fisseler, and F. Weichert, ``Comparative analysis of the accuracy and robustness of the Leap Motion Controller 2,'' \textit{Sensors}, vol. 25, no. 24, p. 7473, 2025, doi: 10.3390/s25247473.

\bibitem{rastgoo2025deep}
R. Rastgoo, K. Kiani, and S. Escalera, ``A deep generative skeleton-based dynamic hand gesture production model,'' \textit{Multimedia Tools and Applications}, vol. 84, no. 40, pp. 48589--48608, 2025, doi: 10.1007/s11042-025-21069-6.

\bibitem{dhanjal2025multilingual}
A. S. Dhanjal and W. Singh, ``Multilingual speech to Indian sign language translation using synthetic animation: a resource-efficient approach,'' \textit{Multimedia Tools and Applications}, vol. 84, no. 21, pp. 24637--24669, 2025, doi: 10.1007/s11042-024-20520-4.

\bibitem{casiez2012oneeuro}
G. Casiez, N. Roussel, and D. Vogel, ``$1\text{€}$ filter: a simple speed-based low-pass filter for noisy input in interactive systems,'' in \textit{Proceedings of the SIGCHI Conference on Human Factors in Computing Systems (CHI)}, 2012, pp. 2527--2530.

\end{thebibliography}

\end{document}